\subsection{Spannungsverlauf in einer Brückengleichrichterschaltung}
% Alt: Aufgabe 9
\Aufgabe{
  \label{Aufg9}
  Skizzieren Sie die Spannung $U_1$ der gegebenen Schaltung.
  Beachten Sie, dass es sich bei $u_\mathrm{E}$ um eine Wechselspannung handelt.
  \begin{figure}[H]
    \centering
    \input{Tikz/src/FigDiodeSpannugsverlauf.tex}\alttext{Die Abbildung zeigt eine elektrische Schaltung mit vier Dioden, die in Form einer Brücke angeordnet sind. Die Dioden sind als Dreiecke mit einem Strich an der Spitze dargestellt, die jeweils in unterschiedliche Richtungen zeigen. Die vier Dioden bilden ein Rechteck. An zwei gegenüberliegenden Ecken des Rechtecks sind zwei Anschlüsse für eine Spannungsquelle angebracht. An den beiden anderen gegenüberliegenden Ecken befinden sich zwei Ausgangsanschlüsse. Die Spannungsquelle ist gezeichnet mit \(u_E\) und der Ausgang der Schaltung mit \(U_1\)}
    \label{fig:FigDiodeSpannugsverlauf}
  \end{figure}
}


\Loesung{
  \begin{itemize}
    \item Bei einer negativen Halbwelle fließt der Strom über $D_{\text{2}}$ und $D_{\text{3}}$.
  \end{itemize}
  \begin{figure}[H]
    \centering
    \input{Tikz/src/LsgDiodeSpannungsverlauf1}
    \label{fig:LsgDiodeSpannungsverlauf1}
  \end{figure}

  \begin{itemize}
    \item Bei einer positiven Halbwelle fließt der Strom über $D_{\text{4}}$ und $D_{\text{1}}$.
  \end{itemize}
  \begin{figure}[H]
    \centering
    \input{Tikz/src/LsgDiodeSpannungsverlauf2}
    \label{fig:LsgDiodeSpannungsverlauf2}
  \end{figure}

  Siehe: Abschnitt \ref{sec:Brückengleichrichter}
}